\documentclass[11pt]{article}
\usepackage[utf8]{inputenc}
\usepackage[spanish, es-tabla]{babel}
\usepackage{amsmath, amssymb}
\usepackage{booktabs}
\usepackage{enumitem}
\usepackage[margin=2.5cm]{geometry}

\setlength{\parindent}{0pt}
\setlength{\parskip}{0.65em}

\title{Econometría ICS-294\\Guía de ejercicios: Clase 8}
\author{Francisco Alfaro Medina}
\date{}

\begin{document}
\maketitle

\section*{Objetivo}

Practicar la estimación de la varianza del error, la matriz de varianzas-covarianzas de MCO, los errores estándar y la intuición detrás de la precisión de los estimadores.

\section*{Ejercicios}

\begin{enumerate}[label=\textbf{Ejercicio \arabic*.}, leftmargin=*]

\item \textbf{Estimador de \(\sigma^2\).}

En una regresión con \(n=50\), \(k=3\) y \(SSR=920\):

\begin{enumerate}[label=\alph*)]
    \item Calcule \(\hat{\sigma}^2\).
    \item Explique por qué se divide por \(n-k-1\) y no por \(n\).
    \item Calcule \(\hat{\sigma}\).
\end{enumerate}

\item \textbf{Matriz de varianzas.}

Suponga que:
\[
(X^\top X)^{-1}=
\begin{bmatrix}
0.20 & -0.04\\
-0.04 & 0.05
\end{bmatrix},
\qquad
\hat{\sigma}^2=100.
\]

\begin{enumerate}[label=\alph*)]
    \item Calcule \(\hat{V}(\hat{\beta})\).
    \item Identifique la varianza de \(\hat{\beta}_1\).
    \item Calcule \(se(\hat{\beta}_1)\).
\end{enumerate}

\item \textbf{Regresión simple.}

En el modelo \(y=\beta_0+\beta_1x+u\), se tiene:
\[
\widehat{Var}(\hat{\beta}_1)=\frac{\hat{\sigma}^2}{\sum_i(x_i-\bar{x})^2}.
\]

\begin{enumerate}[label=\alph*)]
    \item Calcule \(\widehat{Var}(\hat{\beta}_1)\) si \(\hat{\sigma}^2=36\) y \(\sum_i(x_i-\bar{x})^2=144\).
    \item Calcule el error estándar.
    \item ¿Qué pasa si \(\sum_i(x_i-\bar{x})^2\) se duplica?
\end{enumerate}

\item \textbf{Intuición de precisión.}

Explique cómo cambia la precisión de \(\hat{\beta}_1\) en cada caso:

\begin{enumerate}[label=\alph*)]
    \item aumenta el tamaño muestral;
    \item aumenta el ruido no explicado en \(y\);
    \item \(x_1\) casi no varía en la muestra;
    \item \(x_1\) está muy correlacionada con otros regresores.
\end{enumerate}

\item \textbf{Multicolinealidad imperfecta.}

Considere:
\[
Var(\hat{\beta}_1\mid X)=\frac{\sigma^2}{SST_1(1-R_1^2)}.
\]

\begin{enumerate}[label=\alph*)]
    \item Explique qué mide \(R_1^2\).
    \item ¿Qué ocurre si \(R_1^2\) se acerca a 1?
    \item ¿Por qué esto no necesariamente sesga \(\hat{\beta}_1\), pero sí lo vuelve impreciso?
\end{enumerate}

\item \textbf{Cálculo con \(R_1^2\).}

Suponga que \(\sigma^2=25\), \(SST_1=100\). Calcule \(Var(\hat{\beta}_1\mid X)\) para:

\begin{enumerate}[label=\alph*)]
    \item \(R_1^2=0\);
    \item \(R_1^2=0.5\);
    \item \(R_1^2=0.9\).
\end{enumerate}

Interprete la relación entre colinealidad y precisión.

\item \textbf{Errores estándar y significancia.}

Dos estimaciones entregan el mismo coeficiente \(\hat{\beta}_1=2\), pero distintos errores estándar:

\begin{center}
\begin{tabular}{lcc}
\toprule
Modelo & \(\hat{\beta}_1\) & \(se(\hat{\beta}_1)\) \\
\midrule
A & 2 & 0.5 \\
B & 2 & 2.0 \\
\bottomrule
\end{tabular}
\end{center}

\begin{enumerate}[label=\alph*)]
    \item Calcule el estadístico \(t\) contra \(H_0:\beta_1=0\) en cada modelo.
    \item ¿Cuál entrega evidencia estadística más fuerte?
    \item Explique la diferencia.
\end{enumerate}

\item \textbf{Trade-off entre sesgo y varianza.}

Al estudiar el efecto de educación sobre salario, se duda si incluir habilidad medida con ruido.

\begin{enumerate}[label=\alph*)]
    \item ¿Por qué incluirla podría reducir sesgo?
    \item ¿Por qué incluirla podría aumentar varianza?
    \item Explique el trade-off en palabras.
\end{enumerate}

\end{enumerate}

\section*{Preguntas de cierre}

\begin{enumerate}[label=\arabic*.]
    \item ¿Qué es un error estándar?
    \item ¿Qué factores hacen que un estimador sea más preciso?
    \item ¿Por qué la multicolinealidad afecta la varianza?
    \item ¿Cuál es la diferencia entre \(\sigma^2\) y \(\hat{\sigma}^2\)?
\end{enumerate}

\end{document}
